\begin{HVTPage}{Capítulo 6}{Arquitectura modular}{Integración por bloques funcionales}
\HVTLead{Una planta multirruta no debe operar como una sola caja de proceso. Cada módulo conserva sus protecciones y controles, mientras comparte servicios auxiliares cuando sea técnica y operativamente viable.}
\begin{HVTGrid}
\begin{HVTColumn}{Módulo A · Electrólisis}Agua tratada, potencia, electrolizador, separación y secado.\end{HVTColumn}
\begin{HVTColumn}{Módulo B · Pirólisis}Preparación de residuos, reactor, separación de sólidos y tratamiento del gas.\end{HVTColumn}
\begin{HVTColumn}{Módulo C · Biohidrógeno}Pretratamiento, biorreactor, separación gas-líquido y limpieza.\end{HVTColumn}
\begin{HVTColumn}{Módulo D · Servicios comunes}Analítica, purificación final, compresión, almacenamiento y control.\end{HVTColumn}
\end{HVTGrid}
\end{HVTPage}

\begin{HVTPage}{Capítulo 6}{Arquitectura modular}{Entregables de ingeniería conceptual}
\begin{HVTTask}{Entregables de ingeniería conceptual}\item Diagrama de bloques y PFD por ruta. \item Balances preliminares de masa y energía. \item Lista de equipos, instrumentos y servicios. \item Filosofía de operación y parada. \item Matriz de interfaces entre módulos.\end{HVTTask}
\end{HVTPage}
